% ============================================================
% boma_ii.tex
% BOMA Paper II:
% "Compliance, Correction, and Convergence:
%  Operational Specifications for the Branch-Oriented
%  Mathematical Architecture"
% ============================================================

\documentclass[12pt, a4paper]{article}

% ==============================
% PACKAGES
% ==============================
\usepackage[utf8]{inputenc}
\usepackage[T1]{fontenc}
\usepackage{lmodern}
\usepackage{microtype}
\usepackage{amsmath}
\usepackage{amssymb}
\usepackage{amsthm}
\usepackage{mathtools}
\usepackage{xcolor}
\usepackage[
  colorlinks=true,
  linkcolor=blue!60!black,
  citecolor=blue!60!black,
  urlcolor=blue!60!black
]{hyperref}
\usepackage{cleveref}
\usepackage[
  top=2.5cm, bottom=2.5cm,
  left=3cm,  right=3cm
]{geometry}
\usepackage{setspace}
\onehalfspacing
\usepackage{booktabs}
\usepackage{array}
\usepackage{listings}
\usepackage[
  backend=biber,
  style=numeric,
  sorting=nyt,
  maxbibnames=99
]{biblatex}
\addbibresource{boma_ii.bib}

% ==============================
% LISTINGS (YAML)
% ==============================
\lstdefinelanguage{yaml}{
  keywords={ID, Title, Status, EpistemicStatus, EpistemicJustification,
            Conservativity, ConservativityBasis, Depends_on, Introduces,
            Exports, Injects, Derives, DecisionPoint, Options, Selection,
            SelectionRationale, IndependenceBasis, UnChosenOptions,
            ErratumStatus, Nature, CorrectionMode, CorrectionRef,
            FormalPayload, Language, Source, VerificationStatus,
            Domain, Constants, Relations, Functions, Verification,
            Type, Model, Date, Location, Description, Witness,
            Prescription},
  sensitive=false,
  comment=[l]{\#},
  morestring=[b]",
  morestring=[b]'
}
\lstset{
  language=yaml,
  basicstyle=\ttfamily\small,
  keywordstyle=\bfseries,
  commentstyle=\itshape\color{gray},
  stringstyle=\color{black},
  frame=single,
  breaklines=true,
  columns=fullflexible,
  keepspaces=true,
  xleftmargin=1em,
  xrightmargin=1em
}

% ==============================
% THEOREM ENVIRONMENTS
% ==============================
\theoremstyle{plain}
\newtheorem{theorem}{Theorem}[section]
\newtheorem{proposition}[theorem]{Proposition}
\newtheorem{lemma}[theorem]{Lemma}
\newtheorem{corollary}[theorem]{Corollary}

\theoremstyle{definition}
\newtheorem{definition}[theorem]{Definition}
\newtheorem{axiom}[theorem]{Axiom}
\newtheorem{principle}[theorem]{Principle}
\newtheorem{specification}[theorem]{Specification}

\theoremstyle{remark}
\newtheorem{remark}[theorem]{Remark}
\newtheorem{observation}[theorem]{Observation}
\newtheorem{example}[theorem]{Example}

% ==============================
% CUSTOM COMMANDS
% ==============================
\newcommand{\BOMA}{\textsf{BOMA}}
\newcommand{\FS}[1]{\mathfrak{F}(#1)}
\newcommand{\Th}[1]{\mathrm{Th}(#1)}
\newcommand{\Inj}[1]{\mathrm{Inj}(#1)}
\newcommand{\Der}[1]{\mathrm{Der}(#1)}
\newcommand{\IW}[1]{\omega(#1)}
\newcommand{\CIW}[1]{\Omega(#1)}
\newcommand{\DC}[1]{\overline{\mathrm{Dep}}(#1)}
\newcommand{\CC}[1]{\widehat{\mathrm{Content}}(#1)}
\newcommand{\Sep}[1]{\mathrm{Sep}(#1)}
\newcommand{\Inv}[1]{\mathrm{Inv}(#1)}
\newcommand{\Conv}[1]{\mathrm{Conv}(#1)}
\newcommand{\ES}[1]{\mathrm{ES}(#1)}
\newcommand{\Cons}[1]{\mathrm{Cons}(#1)}
\newcommand{\BI}[1]{\mathrm{BI}(#1)}
\newcommand{\CR}[1]{\mathrm{CR}(#1)}
\newcommand{\CD}[3]{\mathrm{CD}(#1,#2,#3)}
\newcommand{\Prov}[2]{\mathrm{Prov}(#1,#2)}
\newcommand{\LN}{\mathrm{LN}}
\newcommand{\FN}{\mathrm{FN}}
\newcommand{\AN}{\mathrm{AN}}
\newcommand{\MC}{\mathrm{MC}}
\newcommand{\C}{\mathrm{C}}
\newcommand{\NC}{\mathrm{NC}}

% ==============================
% TITLE
% ==============================
\title{%
  \textbf{Compliance, Correction, and Convergence}\\[0.5em]
  \large Operational Specifications for the\\
  Branch-Oriented Mathematical Architecture%
}

\author{%
  [Author Name]\\
  \small [Institution]\\
  \small \texttt{[email]}%
}

\date{\today}

% ============================================================
\begin{document}
% ============================================================

\maketitle

% ==============================
% ABSTRACT
% ==============================
\begin{abstract}
The Branch-Oriented Mathematical Architecture (\BOMA{}), introduced
in Paper~I of this series, proposes an organizational paradigm for
mathematical foundations in which every development is structured as a
Branch Tree of locally consistent paths, every genuine choice is
documented at an explicit Decision Point, and every new assumption is
individually declared. Paper~I established the framework's theoretical
architecture but left open a central operational question: what does
it mean, precisely, for a mathematical development to be
\emph{compliant} with \BOMA{}, and how does compliance manifest in
practice?

This paper answers that question in four contributions. First, we
formalize the notion of a \emph{\BOMA{}-compliant development} through
six architectural properties --- Traceability, Transparency,
Auditability, Minimal-Brick Construction, Branchability, and
Correctability --- and establish their formal relationships to the
structural components of Paper~I.

Second, we enrich the binary classification of Paper~I (Forced Step /
Genuine Choice) into a four-way \emph{Epistemic Status} taxonomy:
Logical Necessity, Foundational Necessity, Architectural Necessity,
and Methodological Choice. This taxonomy makes the Branchability
property precise: only Methodological Choice Blocks generate Decision
Points and potential Branches.

Third, we introduce two structural tools absent from Paper~I. The
\emph{Conservativity Declaration} is a mandatory field in every
Block's structural record, making the distinction between zero-injection
and non-zero-injection Blocks immediately legible. The \emph{Internal
Counter-Model} is an architectural diagnostic that characterizes the
structural capacity of a single path from within, complementing the
inter-path role of separation formulas in Paper~I, and serves as a
structural prescription identifying precisely the injections needed for
an Additive Correction.

Fourth, we distinguish two modes of correction: \emph{Additive
Correction}, appropriate when a development falls short of a claimed
consequence, and \emph{Branch Correction}, appropriate when a
foundational choice requires revision. We demonstrate that the
applicable mode is determined by the flaw type and the Epistemic Status
of the affected Block.

All contributions are demonstrated through a worked example: the
Pre-Numeric Kernel of the Transparent Mathematics project, an active
constructive mathematics development that independently instantiates
\BOMA{}-compliance principles and provides the first empirical test of
the framework's structural predictions.
\end{abstract}

\medskip
\noindent\textbf{MSC 2020:}
03A05, 03B22, 03B70, 00A30, 03F65.

\medskip
\noindent\textbf{Keywords:}
\BOMA{} compliance,
epistemic status,
conservativity,
convergent blocks,
architectural DAG,
internal counter-model,
additive correction,
branch correction,
transparent mathematics.

\tableofcontents
\newpage

% ============================================================
\section{Introduction}
\label{sec:introduction}
% ============================================================

\subsection{What Paper~I Established and What Remains Open}
\label{subsec:paper1}

The central contribution of \BOMA{} Paper~I~\cite{BomaI} was
architectural: it showed that the organizational paradigm of existing
mathematical foundations --- linear presentation of axioms and
theorems, with foundational choices invisible, alternatives suppressed,
and injection structure undeclared --- is one paradigm among possible
paradigms, and proposed a precise alternative. That alternative
organized foundational development as a Branch Tree of locally
consistent paths, required every genuine choice to be documented at an
explicit Decision Point, and defined Architectural Transparency as the
conjunctive property of four legibility conditions.

Paper~I established this framework at the level of mathematical
precision appropriate for a foundational proposal. What it did not
establish was its operational content: the specific conditions a
development must satisfy to be \BOMA{}-compliant, the tools available
for diagnosing structural gaps in an ongoing development, and the
procedure for correcting discovered flaws without violating the
immutability of completed Blocks.

These omissions were deliberate. Paper~I was a framework paper; it
defined the architecture without prescribing a protocol. The protocol
is the subject of the present paper.

\subsection{The Central Question}
\label{subsec:question}

The question this paper addresses is simple to state and non-trivial
to answer precisely:
\begin{quote}
\emph{Given a mathematical development in progress, what conditions
must it satisfy to be \BOMA{}-compliant, and how are violations
corrected?}
\end{quote}

The simplicity is deceptive. Compliance is not a single condition but
a conjunction of six, and the six conditions are not independent: the
Minimal-Brick Construction property (P4) determines which Blocks
generate Decision Points, which in turn determines the Branchability
index (P5), which in turn determines where Branch Correction (P6) is
applicable.

The correction question is equally non-trivial. \BOMA{}-compliant
Blocks are immutable once fixed; a development does not revise its
past. When a flaw is discovered --- as it inevitably will be in any
serious mathematical project --- the correction must be effected
without in-place modification. This paper shows that two
fundamentally different correction modes exist, that they apply in
different circumstances, and that distinguishing them correctly is
essential for maintaining the architectural integrity of the
development.

\subsection{The Role of the Practical Example}
\label{subsec:example}

Alongside the formal development, this paper draws on a worked
example: the Transparent Mathematics project, an active constructive
mathematics initiative building mathematical foundations atomically
from minimal traceable Bricks, with machine-verified formal content
and explicit epistemic declarations~\cite{TransparentMath}.

The project was not designed as a \BOMA{} implementation. It arose
independently, from the same motivating concerns --- opacity of
foundational decisions, invisibility of alternatives, difficulty of
tracing consequences --- and converged on structural choices that
instantiate \BOMA{}-compliance principles in practice. The convergence
is not coincidental: it reflects the fact that the structural
requirements of transparent foundational construction, approached
seriously from any direction, generate similar organizational
discipline.

The project's \texttt{ERRATA} document records a discovered flaw ---
an over-claim about the generative capacity of the Pre-Numeric Kernel
--- together with the counter-model that established it and the
structural prescription for correction. This episode provides the
first empirical test of the framework's structural predictions, and
its analysis motivates the distinction between Additive and Branch
Correction developed in \S\ref{sec:correction} and the formalization
of the Internal Counter-Model in \S\ref{sec:counter}.

\subsection{Contributions}
\label{subsec:contributions}

This paper makes the following contributions.

\textbf{Formalization of \BOMA{}-Compliance.} We define a
\BOMA{}-compliant development through six architectural properties
(\S\ref{sec:compliance}) and establish their formal relationships to
the structural components of Paper~I.

\textbf{Four-Way Epistemic Status.} We replace the binary Forced Step
/ Genuine Choice classification of Paper~I with a four-way taxonomy
--- Logical Necessity, Foundational Necessity, Architectural
Necessity, Methodological Choice --- and show that the distinction
carries structural consequences for the Branchability property
(\S\ref{sec:epistemic}).

\textbf{Conservativity as a Mandatory Field.} We formalize
Conservative and Non-Conservative Block extensions as a mandatory
declaration in every Block's structural record, and show its
relationship to the injection weight of Paper~I (\S\ref{sec:cons}).

\textbf{Two Modes of Correction.} We distinguish Additive Correction
from Branch Correction, characterize the conditions under which each
applies, and prove that the applicable mode is determined by the
Epistemic Status of the Block in which the flaw originates
(\S\ref{sec:correction}).

\textbf{The Internal Counter-Model.} We formalize the Counter-Model
as an internal architectural diagnostic and show that it serves as a
structural prescription identifying precisely the injections needed
for the next Additive Correction (\S\ref{sec:counter}).

\textbf{The Architectural DAG and Convergent Blocks.} We extend the
Branch Tree of Paper~I to the full Block-level Architectural DAG and
introduce the Convergent Block --- a Block appearing identically in
two diverging paths --- as the architectural criterion for
foundational invariance (\S\ref{sec:dag}).

\textbf{The Machine-Readable Block Record.} We formalize the Block
Summary Record as the infrastructure for mechanical Auditability
(\S\ref{sec:record}).

\textbf{Worked Example.} We analyze the Transparent Mathematics
project's Pre-Numeric Kernel as a complete instance of
\BOMA{}-compliance in practice (\S\ref{sec:example}).

\subsection{What This Paper Does Not Claim}
\label{subsec:notclaim}

This paper does not claim that the Transparent Mathematics project was
designed as a \BOMA{} implementation. It was not. The project is
analyzed here as an independently arising instantiation of structural
principles. The convergence between the project's organizational
discipline and \BOMA{}'s structural requirements is evidence of the
framework's descriptive adequacy, not of the project's theoretical
dependence on it.

This paper does not claim that the six compliance properties are
exhaustive or uniquely correct. They represent the properties that
emerge from taking seriously the structural requirements of transparent
foundational construction as identified in Paper~I.

\subsection{Paper Organization}
\label{subsec:organization}

Section~\ref{sec:epistemic} develops the four-way Epistemic Status
taxonomy. Section~\ref{sec:cons} formalizes Conservativity as a
mandatory declaration. Section~\ref{sec:compliance} defines
\BOMA{}-compliance through six architectural properties.
Section~\ref{sec:dag} introduces the Architectural DAG and the
Convergent Block. Section~\ref{sec:correction} distinguishes Additive
from Branch Correction. Section~\ref{sec:counter} formalizes the
Internal Counter-Model. Section~\ref{sec:record} establishes the
Machine-Readable Block Record. Section~\ref{sec:example} analyzes the
Transparent Mathematics project. Section~\ref{sec:paper1} relates the
contributions to Paper~I. Section~\ref{sec:conclusion} concludes.

% ============================================================
\section{Epistemic Status Taxonomy}
\label{sec:epistemic}
% ============================================================

\subsection{Motivation}
\label{subsec:ep-motivation}

Paper~I distinguished between two kinds of steps in an architectural
development: \emph{Forced Steps}, in which the content of the next
Block is uniquely determined by prior cumulative content, and
\emph{Genuine Choices}, in which at least two distinct, logically
admissible, and architecturally non-equivalent continuations are
available. This binary classification carries the essential structural
load of Paper~I.

The binary classification is adequate for Paper~I's purposes, but it
collapses distinctions that are epistemically significant. A step
forced by the underlying logic is different in character from a step
forced by the declared goals of the development: the former cannot be
otherwise under any mathematical circumstances; the latter could be
otherwise if the goal were revised. A step forced by the structural
discipline of the framework is different again. The present section
replaces the binary classification with a four-way \emph{Epistemic
Status} taxonomy that preserves the structural consequences of
Paper~I while making the source of necessity explicit in every case.

\subsection{The Four Epistemic Statuses}
\label{subsec:ep-four}

\begin{definition}[Epistemic Status]
\label{def:es}
The \emph{Epistemic Status} of a Block $\beta$ is a mandatory
declaration drawn from the following four-element set:
\[
  \ES{\beta}
  \in \{\LN,\, \FN,\, \AN,\, \MC\}
\]
where the four values are defined in
Definitions~\ref{def:ln}--\ref{def:mc} below.
\end{definition}

\begin{definition}[Logical Necessity --- $\LN$]
\label{def:ln}
A Block $\beta$ has Epistemic Status \emph{Logical Necessity} if its
content is uniquely forced by the cumulative content of all ancestors
under the declared object logic:
\[
  \ES{\beta} = \LN
  \iff
  \forall \varphi \in \mathrm{Content}(\beta) :
  \CC{\DC{\beta}} \vdash \varphi.
\]
An $\LN$ Block has empty injection set ($\Inj{\beta} = \emptyset$)
and is therefore Conservative.
\end{definition}

\begin{definition}[Foundational Necessity --- $\FN$]
\label{def:fn}
A Block $\beta$ has Epistemic Status \emph{Foundational Necessity} if
its content is not logically forced by prior content alone, but is
forced by the conjunction of prior content and the declared goals of
the development. An $\FN$ Block has non-empty injection set and
is therefore Non-Conservative.
\end{definition}

\begin{remark}
Foundational Necessity makes explicit a form of dependency often left
implicit: the dependency of individual steps on the overall direction
of the development. When a Block is declared $\FN$, the declaration
carries a commitment: the development's goal is declared, and the
Block's relationship to that goal is stated. A reader who disputes
the Block's content is directed to examine the goal declaration rather
than the Block in isolation.
\end{remark}

\begin{definition}[Architectural Necessity --- $\AN$]
\label{def:an}
A Block $\beta$ has Epistemic Status \emph{Architectural Necessity} if
its content is required by the structural discipline of the \BOMA{}
framework itself --- by the architectural paradigm --- rather than by
the mathematical logic or the development's mathematical goals.
$\AN$ Blocks constitute the constitutional layer of a
\BOMA{}-compliant development.
\end{definition}

\begin{definition}[Methodological Choice --- $\MC$]
\label{def:mc}
A Block $\beta$ has Epistemic Status \emph{Methodological Choice} if
its content represents a genuine choice among at least two distinct,
logically admissible, and architecturally non-equivalent alternatives,
none of which is forced by prior cumulative content, the development's
declared goals, or the architectural paradigm. An $\MC$ Block
corresponds precisely to a Genuine Choice in the sense of
Paper~I, Definition~4.3.10, and generates a Decision Point.
\end{definition}

\begin{remark}
The identification of a Block as $\MC$ may be \emph{objective}
(supported by a formal independence result) or \emph{interpretive}
(resting on the developer's judgment), as established in Paper~I,
Remark~4.3.11. In both cases, the basis of identification must be
declared in the Block's Rationale component.
\end{remark}

\subsection{Structural Consequences}
\label{subsec:ep-structure}

\begin{proposition}[Injection Sets and Epistemic Status]
\label{prop:es-inj}
The following relationships hold:
\begin{align*}
  \ES{\beta} = \LN &\implies \Inj{\beta} = \emptyset \\
  \ES{\beta} \in \{\FN,\,\MC\} &\implies \Inj{\beta} \neq \emptyset \\
  \ES{\beta} = \AN &\implies \Inj{\beta} = \emptyset
                      \;\text{ or }\; \Inj{\beta} \neq \emptyset
\end{align*}
\end{proposition}

\begin{proof}
$\LN$ requires all content derivable from prior content, hence empty
injection set. $\FN$ and $\MC$ require content not derivable from
prior content alone, hence non-empty injection set. $\AN$ Blocks
may or may not inject into the object-level theory depending on
whether the architectural discipline is formalized within the same
logical framework.
\end{proof}

\begin{proposition}[Decision Points and Epistemic Status]
\label{prop:es-dp}
A Block $\beta$ is associated with a Decision Point if and only if
$\ES{\beta} = \MC$.
\end{proposition}

\begin{proof}
By Definition~\ref{def:mc}, an $\MC$ Block corresponds to a Genuine
Choice, and by Paper~I, Definition~4.3.10, every Genuine Choice gives
rise to a Decision Point. Blocks with $\ES{\beta} \in \{\LN,\FN,\AN\}$
admit no genuine alternatives within their respective constraints.
\end{proof}

\begin{corollary}[Branchability Index]
\label{cor:bi}
A \BOMA{}-compliant development $\mathcal{P}$ is branchable at Block
$\beta$ if and only if $\ES{\beta} = \MC$. The \emph{Branchability
Index} of $\mathcal{P}$ is:
\[
  \BI{\mathcal{P}}
  = |\{\beta \in \mathcal{P} \mid \ES{\beta} = \MC\}|.
\]
\end{corollary}

\subsection{The Mandatory Epistemic Status Declaration}
\label{subsec:ep-declaration}

\begin{axiom}[Mandatory Epistemic Status Declaration]
\label{ax:es}
Every Block $\beta$ in a \BOMA{}-compliant development must include
an explicit declaration of $\ES{\beta}$, providing:
\begin{enumerate}
  \item one of the four Epistemic Statuses of
        Definition~\ref{def:es};
  \item a justification consistent with the chosen status;
  \item the declaration at the time of Block introduction,
        not retroactively.
\end{enumerate}
\end{axiom}

\begin{proposition}[Compatibility with Paper~I]
\label{prop:es-compat}
The four-way Epistemic Status taxonomy is a refinement of Paper~I's
binary classification:
\begin{align*}
  \text{Forced Step (Paper I)} &\equiv \ES{\beta}
      \in \{\LN,\,\FN,\,\AN\} \\
  \text{Genuine Choice (Paper I)} &\equiv \ES{\beta} = \MC.
\end{align*}
All structural results of Paper~I remain valid under the four-way
taxonomy.
\end{proposition}

% ============================================================
\section{Conservativity as a Mandatory Field}
\label{sec:cons}
% ============================================================

\subsection{Formal Definitions}
\label{subsec:cons-defs}

\begin{definition}[Conservative Extension]
\label{def:conservative}
A Block $\beta$ is a \emph{conservative extension} of the development
up to $\DC{\beta}$ if every statement in its content is derivable
from the cumulative content of its ancestors:
\[
  \mathrm{IsConservative}(\beta)
  \iff
  \Inj{\beta} = \emptyset.
\]
A conservative Block adds no new foundational commitment.
\end{definition}

\begin{definition}[Non-Conservative Extension]
\label{def:nonconservative}
A Block $\beta$ is a \emph{non-conservative extension} if
$\Inj{\beta} \neq \emptyset$. It introduces at least one statement
not derivable from prior content and restricts the class of
admissible models.
\end{definition}

\begin{definition}[Conservativity Field]
\label{def:cons-field}
The \emph{Conservativity Field} of a Block $\beta$ is the binary
declaration:
\[
  \Cons{\beta} \in \{\C,\, \NC\}
\]
where $\C$ denotes Conservative and $\NC$ denotes Non-Conservative.
\end{definition}

\subsection{Mandatory Declaration and Relationships}
\label{subsec:cons-mandatory}

\begin{axiom}[Mandatory Conservativity Declaration]
\label{ax:cons}
Every Block $\beta$ in a \BOMA{}-compliant development must include
an explicit declaration of $\Cons{\beta}$ satisfying:
\begin{enumerate}
  \item \textbf{Accuracy:} the declared value matches the injection
        set;
  \item \textbf{Contemporaneity:} the declaration is made at Block
        introduction;
  \item \textbf{Verifiability:} supporting information for
        independent verification is provided.
\end{enumerate}
\end{axiom}

\begin{proposition}[Conservativity and Injection Weight]
\label{prop:cons-iw}
For any Block $\beta$:
\[
  \Cons{\beta} = \C \iff \IW{\beta} = 0 \qquad
  \Cons{\beta} = \NC \iff \IW{\beta} > 0.
\]
\end{proposition}

\begin{proof}
Immediate from Definitions~\ref{def:conservative}--\ref{def:cons-field}
and Paper~I, Definition~4.6.3: injection weight is zero iff the
injection set is empty.
\end{proof}

\begin{proposition}[Conservativity and Epistemic Status]
\label{prop:cons-es}
The following relationships hold:
\begin{align*}
  \ES{\beta} = \LN &\implies \Cons{\beta} = \C \\
  \ES{\beta} \in \{\FN,\,\MC\} &\implies \Cons{\beta} = \NC \\
  \ES{\beta} = \AN &\implies \Cons{\beta} \in \{\C,\,\NC\}.
\end{align*}
\end{proposition}

\begin{corollary}[Conservativity as a Necessary Condition for Branch Origins]
\label{cor:cons-branch}
A Block $\beta$ can be a Branch origin only if $\Cons{\beta} = \NC$.
\end{corollary}

\begin{proof}
By Proposition~\ref{prop:es-dp}, a Branch origin requires
$\ES{\beta} = \MC$. By Proposition~\ref{prop:cons-es},
$\ES{\beta} = \MC$ implies $\Cons{\beta} = \NC$.
\end{proof}

\begin{definition}[Conservativity Ratio]
\label{def:cr}
The \emph{Conservativity Ratio} of a path $\mathcal{P}$ is:
\[
  \CR{\mathcal{P}}
  = \frac{|\{\beta \in \mathcal{P} \mid \Cons{\beta} = \C\}|}
         {|\mathcal{P}|}.
\]
\end{definition}

\begin{proposition}[Conservative Blocks Do Not Restrict Models]
\label{prop:cons-models}
Let $\beta_k$ be a Conservative Block in a locally consistent path.
Then every model satisfying the cumulative content of the preceding
Blocks also satisfies the content of $\beta_k$.
\end{proposition}

\begin{proof}
Since $\Cons{\beta_k} = \C$, every statement in
$\mathrm{Content}(\beta_k)$ is derivable from prior cumulative
content. By soundness of the object logic, any model satisfying the
premises satisfies the conclusion.
\end{proof}

% ============================================================
\section{\BOMA{}-Compliant Development}
\label{sec:compliance}
% ============================================================

\subsection{The Six Properties}
\label{subsec:six}

\begin{definition}[\BOMA{}-Compliant Development]
\label{def:compliant}
A mathematical development $\mathcal{P}$ is \emph{\BOMA{}-compliant}
if it satisfies the following six architectural properties
simultaneously.

\medskip
\noindent\textbf{P1 --- Traceability.}
Every statement $\varphi \in \Th{\mathcal{P}}$ has a structurally
determined provenance:
\begin{align*}
  \Prov{\varphi}{\mathcal{P}}
  = \bigl(&\{\beta_i \mid \varphi \text{ depends on }
    \mathrm{Content}(\beta_i)\},\\
    &\{\mathcal{D}_j \mid \varphi \text{ depends on }
    \mathrm{Selection}(\mathcal{D}_j)\}\bigr)
\end{align*}
recoverable from the structural record without external
reconstruction. In practice, every formal entity carries a Global
Unique Identifier of the form $\mathrm{TYPE\text{-}BLOCK\text{-}SEQ}$.

\medskip
\noindent\textbf{P2 --- Transparency.}
The development is architecturally transparent in the sense of
Paper~I, Definition~4.7.5, and additionally every Block declares
$\ES{\beta}$ (Axiom~\ref{ax:es}) and $\Cons{\beta}$
(Axiom~\ref{ax:cons}).

\medskip
\noindent\textbf{P3 --- Auditability.}
Any researcher with access to the structural record can independently:
(i) verify local consistency of any path;
(ii) recover the provenance of any statement;
(iii) identify the Decision Points and their documented alternatives;
(iv) verify Conservativity declarations.
\emph{Mechanical Auditability} is achieved when the object logic
supports automated proof checking; \emph{Interpretive Auditability}
is the minimum requirement.

\medskip
\noindent\textbf{P4 --- Minimal-Brick Construction.}
Every Block satisfies the Minimal Injection Principle (Paper~I,
Axiom~4.6.6) with two additional disciplines:

\begin{itemize}
  \item \textbf{Vocabulary-Axiom Decoupling:} the introduction of a
        vocabulary item and the introduction of axioms governing it
        occupy separate Blocks.
  \item \textbf{Single-Purpose Commitment:} each Non-Conservative
        Block introduces exactly one logically independent core
        concept.
\end{itemize}

\medskip
\noindent\textbf{P5 --- Branchability.}
The Branchability Index $\BI{\mathcal{P}}$ is accurately computed
and every $\MC$ Block is associated with a fully documented Decision
Point satisfying Paper~I, Axiom~4.3.4. Non-$\MC$ Blocks are not
Branch origins, and the framework imposes no requirement that every
Block be branchable.

\medskip
\noindent\textbf{P6 --- Correctability.}
Every discovered flaw is remedied through one of two architecturally
admissible correction procedures (Additive or Branch Correction, as
defined in \S\ref{sec:correction}), and no completed Block is
modified in place. Completed Blocks are immutable; their
\texttt{ErratumStatus} field is the only field that may be updated
after the Block is marked Immutable.
\end{definition}

\subsection{Formal Definition}
\label{subsec:formal}

\begin{definition}[\BOMA{}-Compliance --- Formal]
\label{def:formal-compliant}
A mathematical development $\mathcal{P}$ is \emph{\BOMA{}-compliant}
if and only if:
\begin{enumerate}
  \item $\mathcal{P}$ is architecturally transparent (Paper~I,
        Definition~4.7.5);
  \item every Block declares $\ES{\beta}$ satisfying
        Axiom~\ref{ax:es};
  \item every Block declares $\Cons{\beta}$ satisfying
        Axiom~\ref{ax:cons};
  \item every Block satisfies P4 (Minimal-Brick Construction);
  \item $\BI{\mathcal{P}}$ is accurately computed and every $\MC$
        Block has a fully documented Decision Point;
  \item every flaw is remedied through an admissible correction
        procedure with no in-place modification.
\end{enumerate}
\end{definition}

\subsection{Compliance Hierarchy and Relations to Paper~I}
\label{subsec:hierarchy}

\begin{proposition}[\BOMA{}-Compliance Strengthens Architectural Transparency]
\label{prop:comp-at}
Every \BOMA{}-compliant development is architecturally transparent
(Paper~I, Definition~4.7.5). The converse does not hold:
Architectural Transparency does not require Epistemic Status,
Conservativity, Voc\-ab\-u\-lary-Axiom Decoupling, or the ERRATA
mechanism.
\end{proposition}

\begin{proposition}[\BOMA{}-Compliance and Local Consistency]
\label{prop:comp-lc}
Every locally consistent path in a \BOMA{}-compliant development
remains locally consistent after any Additive Correction, and after
any Branch Correction that selects a locally consistent alternative.
\end{proposition}

\begin{proof}
Additive Correction introduces new Blocks that are locally consistent
in isolation. No completed Block is modified in place; the original
path's local consistency is unaffected. Branch Correction opens a new
path whose local consistency is assessed independently by Paper~I,
Principle~4.5.5.
\end{proof}

% ============================================================
\section{The Architectural DAG and Convergent Blocks}
\label{sec:dag}
% ============================================================

\subsection{The Architectural DAG}
\label{subsec:dag-def}

\begin{definition}[Architectural DAG]
\label{def:dag}
The \emph{Architectural DAG} of a \BOMA{}-compliant development is a
directed acyclic graph $\mathcal{G} = (V, E)$ where:
\begin{itemize}
  \item $V$ is the set of all Blocks appearing in any locally
        consistent path of the development;
  \item $E$ contains a directed edge $(\beta, \beta')$ iff
        $\beta \in \mathrm{Dep}(\beta')$ for some path.
\end{itemize}
Acyclicity follows from Paper~I, Axiom~4.2.4, together with the
global identifier ordering (Block $\beta \in \mathrm{Dep}(\beta')$
implies $\mathrm{Id}(\beta) < \mathrm{Id}(\beta')$).
\end{definition}

\begin{definition}[Branch Tree as Projection]
\label{def:projection}
The Branch Tree $\mathcal{T}$ of Paper~I, Definition~4.4.12, is the
\emph{Decision-Level Projection} of $\mathcal{G}$: the sub-DAG
induced by Decision Blocks alone, with transitive edges between
consecutive Decision Blocks.
\end{definition}

\begin{remark}
The projection from $\mathcal{G}$ to $\mathcal{T}$ is lossy: it
retains the decision-level structure but discards the Block-level
structure between Decision Points. Two developments with identical
Branch Trees may have very different Architectural DAGs.
\end{remark}

\subsection{Convergent Blocks}
\label{subsec:convergent}

\begin{definition}[Convergent Block]
\label{def:convergent}
A Block $\beta$ is a \emph{Convergent Block} with respect to locally
consistent paths $\mathcal{P}_1$ and $\mathcal{P}_2$ diverging at
Decision Point $\mathcal{D}$ if:
\begin{enumerate}
  \item $\beta \in \mathcal{P}_1$ and $\beta \in \mathcal{P}_2$;
  \item $\mathrm{Content}(\beta)$ is identical in both paths;
  \item $\beta \notin \mathcal{P}_1 \wedge \mathcal{P}_2$ (the Block
        is not part of the common ancestry).
\end{enumerate}
Condition~3 distinguishes Convergent Blocks from ordinary shared
ancestors: a Convergent Block is reached independently after the
divergence at $\mathcal{D}$.
\end{definition}

\begin{definition}[Convergence Point and Depth]
\label{def:conv-point}
A \emph{Convergence Point} in $\mathcal{G}$ is a vertex
$\beta \in \Conv{\mathcal{G}}$ that is a Convergent Block with
respect to at least two paths. The \emph{convergence depth} of
$\beta$ with respect to $\mathcal{P}_1$ and $\mathcal{P}_2$ is:
\[
  \CD{\beta}{\mathcal{P}_1}{\mathcal{P}_2}
  = \min\!\bigl(
      |\mathcal{P}_1 \setminus (\mathcal{P}_1 \wedge \mathcal{P}_2)|,\;
      |\mathcal{P}_2 \setminus (\mathcal{P}_1 \wedge \mathcal{P}_2)|
    \bigr).
\]
\end{definition}

\subsection{Path-Independence and Mathematical Significance}
\label{subsec:path-indep}

\begin{theorem}[Path-Independence of Convergent Blocks]
\label{thm:path-indep}
Let $\beta$ be a Convergent Block with respect to locally consistent
paths $\mathcal{P}_1$ and $\mathcal{P}_2$ diverging at $\mathcal{D}$.
Then:
\begin{enumerate}
  \item The content of $\beta$ is architecturally present and
        structurally valid in both paths.
  \item The content of $\beta$ is \emph{path-independent} with
        respect to $\mathcal{D}$: the selection at $\mathcal{D}$
        was not sufficient to prevent $\beta$'s content from being
        reached in either path.
  \item $\mathrm{Content}(\beta) \subseteq
        \Th{\mathcal{P}_1} \cap \Th{\mathcal{P}_2}$.
  \item $\mathcal{D}$ is non-trivial, and $\beta$'s content is
        precisely what $\mathcal{D}$'s non-triviality does
        \emph{not} affect.
\end{enumerate}
\end{theorem}

\begin{proof}
(1) By Definition~\ref{def:convergent}, $\beta$ appears in both
paths with identical content; whether Conservative or Non-Conservative,
it is architecturally valid in both.

(2) Since $\beta \notin \mathcal{P}_1 \wedge \mathcal{P}_2$, it was
reached after the divergence. Both paths arrived at $\beta$ with
identical content despite different post-Decision histories. The
content is therefore not determined by the selection at $\mathcal{D}$.

(3) $\mathrm{Content}(\beta) \subseteq \Th{\mathcal{P}_1}$ and
$\mathrm{Content}(\beta) \subseteq \Th{\mathcal{P}_2}$ follow from
(1); the intersection claim follows.

(4) Paper~I, Definition~4.5.7 establishes non-triviality from the
divergence of locally consistent paths. The content of $\beta$, being
common to both, falls outside the scope of $\mathcal{D}$'s
non-triviality.
\end{proof}

\begin{definition}[Decision-Point Invariant]
\label{def:dpi}
The \emph{Decision-Point Invariant} of $\mathcal{D}$ is:
\[
  \Inv{\mathcal{D}}
  = \bigl\{\beta \;\big|\; \beta \text{ is a Convergent Block
    w.r.t.\ some }
    \mathcal{P}_1, \mathcal{P}_2 \text{ diverging at }
    \mathcal{D}\bigr\}.
\]
\end{definition}

\begin{proposition}[Complementarity of Invariants and Separation Formulas]
\label{prop:complement}
For any non-trivial $\mathcal{D}$ with paths $\mathcal{P}_1$,
$\mathcal{P}_2$:
\[
  \mathrm{Content}(\beta) \in \Inv{\mathcal{D}}
  \implies
  \mathrm{Content}(\beta) \notin \Sep{\mathcal{D}}.
\]
The content of Convergent Blocks and the separation formulas of
Paper~I are disjoint.
\end{proposition}

\begin{proof}
A separation formula $\varphi \in \Sep{\mathcal{D}}$ satisfies
$\mathcal{P}_1 \vdash \varphi$ and $\mathcal{P}_2 \vdash \neg\varphi$.
If $\varphi \in \mathrm{Content}(\beta)$ for a Convergent Block
$\beta$, then $\mathcal{P}_2 \vdash \varphi$, contradicting
$\mathcal{P}_2 \vdash \neg\varphi$ given local consistency of
$\mathcal{P}_2$.
\end{proof}

\begin{corollary}[Architectural Portrait of a Decision Point]
\label{cor:portrait}
The complete architectural portrait of a non-trivial $\mathcal{D}$
with respect to $\mathcal{P}_1$ and $\mathcal{P}_2$ consists of:
the common ancestry $\mathcal{P}_1 \wedge \mathcal{P}_2$;
the separation formulas $\Sep{\mathcal{D}}$ (what $\mathcal{D}$
makes different);
the Decision-Point Invariant $\Inv{\mathcal{D}}$ (what $\mathcal{D}$
leaves unchanged);
and the injection divergence $\Delta_{\mathrm{Inj}}(\mathcal{P}_1,
\mathcal{P}_2)$ (Paper~I, Definition~6.3.5).
\end{corollary}

% ============================================================
\section{Two Types of Correction}
\label{sec:correction}
% ============================================================

\subsection{The ERRATA Mechanism}
\label{subsec:errata}

\begin{definition}[ERRATA Entry]
\label{def:errata}
An \emph{ERRATA Entry} for a discovered flaw is a formal record:
\[
  \mathrm{ERR}
  = \bigl(\mathrm{Id},\;
          \mathrm{Location},\;
          \mathrm{Nature},\;
          \mathrm{Witness},\;
          \mathrm{Correction}\bigr)
\]
where:
\begin{itemize}\setlength{\itemsep}{0pt}
  \item $\mathrm{Nature} \in \{
        \mathrm{LogicalInconsistency},\;
        \mathrm{OverClaim},$\\
        $\phantom{\mathrm{Nature} \in \{}
        \mathrm{MisclassifiedEpistemic},\;
        \mathrm{ViolatedMinimalInjection}\}$
  \item $\mathrm{Witness}$ is a formal demonstration of the flaw:
        a derivation of $\bot$, a counter-model, or a derivation of
        a supposedly injected statement
  \item $\mathrm{Correction} \in \{\mathrm{Additive},\,
        \mathrm{Branch}\}$
\end{itemize}
\end{definition}

\begin{axiom}[Mandatory ERRATA Logging]
\label{ax:errata}
Every discovered flaw must be formally logged as an ERRATA Entry
before any correction is introduced. The flawed Block receives an
\texttt{ErratumStatus} field pointing to $\mathrm{Id}(\mathrm{ERR})$.
The ERRATA Entry is itself immutable once logged.
\end{axiom}

\subsection{The Two Correction Modes}
\label{subsec:modes}

\begin{definition}[Additive Correction]
\label{def:additive}
An \emph{Additive Correction} introduces new Blocks into the current
path without opening a Branch and without modifying any existing
Block. If the current path is $\mathcal{P} = (\beta_0,\ldots,\beta_n)$,
an Additive Correction produces:
\[
  \mathcal{P}^+
  = (\beta_0, \ldots, \beta_n, \gamma_1, \gamma_2, \ldots, \gamma_k)
\]
where $\gamma_1,\ldots,\gamma_k$ are new Blocks addressing the flaw.
\end{definition}

\begin{definition}[Branch Correction]
\label{def:branch-correction}
A \emph{Branch Correction} opens a new Branch at the Decision Point
associated with the flawed $\MC$ Block, selecting a different option.
If the flawed Block $\beta_j$ corresponds to Decision Point
$\mathcal{D}_j$ with selection $O_s$, a Branch Correction selects
some $O_i \neq O_s$ and opens:
\[
  \mathcal{B}_i
  = (\beta_0, \ldots, \beta_{j-1},\;
     \gamma^i_{j}, \gamma^i_{j+1}, \ldots)
\]
where $(\beta_0,\ldots,\beta_{j-1})$ is the branching prefix
inherited exactly.
\end{definition}

\subsection{Correction Mode Determination}
\label{subsec:mode-det}

\begin{theorem}[Correction Mode Determination]
\label{thm:mode}
The applicable correction mode is determined by the flaw type and the
Epistemic Status of the affected Block as follows:

\medskip
\begin{center}
\begin{tabular}{p{6cm}p{2.2cm}p{2cm}}
\toprule
Flaw Type & $\ES{\beta}$ & Mode \\
\midrule
$\mathrm{OverClaim}$ & Any & Additive \\
$\mathrm{LogicalInconsistency}$ & $\LN$ or $\FN$ & Additive \\
$\mathrm{LogicalInconsistency}$ & $\MC$ & Branch \\
$\mathrm{MisclassifiedEpistemic}$:
  $\LN/\FN/\AN$ declared, $\MC$ correct & Any & Branch \\
$\mathrm{MisclassifiedEpistemic}$:
  $\MC$ declared, $\LN/\FN/\AN$ correct & Any & Additive \\
$\mathrm{ViolatedMinimalInjection}$ & Any & Additive \\
\bottomrule
\end{tabular}
\end{center}
\end{theorem}

\begin{proof}[Proof sketch]
\emph{Over-Claim:} The flaw is in attributed consequences, not in any
Block's content. Existing Blocks are sound; the path must be extended
with the missing injections. No Branch is needed.

\emph{Logical Inconsistency, $\LN$/$\FN$:} Content was forced by
logic or goals. Inconsistency traces to earlier injections that
combine to produce a contradiction. Corrective Blocks address the
source additively.

\emph{Logical Inconsistency, $\MC$:} The chosen option at a Decision
Point led to inconsistency. A different option must be selected:
Branch Correction opens the alternative path.

\emph{Misclassified $\LN/\FN/\AN$ when $\MC$ is correct:} A genuine
alternative was suppressed. The correction opens the Decision Point
that was previously invisible: Branch Correction applies.

\emph{Misclassified $\MC$ when $\LN/\FN/\AN$ is correct:} A spurious
Decision Point was declared. A supplementary Additive Block formally
establishes that the content was forced.

\emph{Violated Minimal Injection:} A Block injected more than
necessary. A corrective Additive Block replaces the strong injection
with the weaker sufficient statement.
\end{proof}

\begin{proposition}[Incoherence of Wrong-Mode Correction]
\label{prop:wrong-mode}
Applying Branch Correction to a flaw that requires Additive
Correction, or vice versa, produces architecturally incoherent
results.
\end{proposition}

\begin{proof}
Branch Correction applied to an Over-Claim changes the foundational
option selected but does not supply the missing injections; the
over-claim persists. Additive Correction applied to a Logical
Inconsistency in an $\MC$ Block extends an already-invalid path
(by explosion) without resolving the inconsistency at its source.
In each case the flaw is not addressed by the applied mode.
\end{proof}

% ============================================================
\section{The Internal Counter-Model}
\label{sec:counter}
% ============================================================

\subsection{Formal Definitions}
\label{subsec:cm-defs}

\begin{definition}[Internal Counter-Model]
\label{def:cm}
Let $\mathcal{P}$ be a locally consistent \BOMA{}-compliant path and
$\varphi$ a formula in the path's declared object logic. An
\emph{internal counter-model} for $\varphi$ relative to $\mathcal{P}$
is a model $\mathcal{M}$ satisfying:
\begin{enumerate}
  \item \emph{Compatibility:} $\mathcal{M} \models
        \CC{\mathcal{P}}$;
  \item \emph{Falsification:} $\mathcal{M} \not\models \varphi$.
\end{enumerate}
The existence of $\mathcal{M}$ establishes $\mathcal{P} \nvdash
\varphi$.
\end{definition}

\subsection{Existence and Non-Derivability}
\label{subsec:cm-existence}

\begin{theorem}[Counter-Model Existence]
\label{thm:cm-exist}
Let $\mathcal{P}$ be locally consistent. For any formula $\varphi$:
\[
  \mathcal{P} \nvdash \varphi
  \iff
  \exists\,\mathcal{M} :
  \mathcal{M} \models \CC{\mathcal{P}}
  \;\land\;
  \mathcal{M} \not\models \varphi.
\]
\end{theorem}

\begin{proof}
Right-to-left: soundness of the object logic. Left-to-right: for
logics with a completeness theorem (classical first-order logic via
Gödel's completeness theorem~\cite{Enderton2001}; intuitionistic
logic via Kripke completeness~\cite{Heyting1956}), if
$\mathcal{P} \nvdash \varphi$ then $\CC{\mathcal{P}} \cup \{\neg\varphi\}$
is consistent and has a model $\mathcal{M}$.
\end{proof}

\subsection{Counter-Models as Architectural Diagnostics}
\label{subsec:cm-diagnostic}

\begin{theorem}[Counter-Model as Structural Prescription]
\label{thm:cm-prescribe}
Let $\mathcal{M}$ be an internal counter-model for $\varphi$ relative
to locally consistent $\mathcal{P}$. Then:
\begin{enumerate}
  \item \textbf{Consistency confirmation:} $\mathcal{P}$ is locally
        consistent.
  \item \textbf{Structural deficit identification:} The conditions
        satisfied by $\mathcal{M}$ but incompatible with $\varphi$
        constitute the precise structural deficit of $\mathcal{P}$
        with respect to $\varphi$.
  \item \textbf{Additive Correction prescription:} The minimal set
        of injections that would rule out $\mathcal{M}$ constitutes
        the minimal Additive Correction required to make $\varphi$
        derivable.
  \item \textbf{Non-branching character:} The counter-model does not
        compare $\mathcal{P}$ to any other path; it characterizes
        $\mathcal{P}$ from within.
\end{enumerate}
\end{theorem}

\begin{proof}
(1) $\mathcal{M} \models \CC{\mathcal{P}}$ establishes satisfiability,
hence consistency.
(2) The structural deficit $S(\mathcal{M},\varphi)$ consists of
conditions in $\mathcal{M}$ not ruled out by $\CC{\mathcal{P}}$ that
permit $\varphi$ to be falsified.
(3) The negations of elements of $S(\mathcal{M},\varphi) \setminus
\CC{\mathcal{P}}$ are the required injections; minimality follows from
the maximally informative counter-model construction.
(4) The construction requires only $\mathcal{P}$ and $\varphi$; no
other path is referenced.
\end{proof}

\begin{proposition}[Counter-Model Confirms Local Consistency]
\label{prop:cm-lc}
If an internal counter-model for any $\varphi$ relative to
$\mathcal{P}$ exists, then $\mathcal{P}$ is locally consistent.
\end{proposition}

\begin{proof}
$\mathcal{M} \models \CC{\mathcal{P}}$ implies satisfiability of
$\CC{\mathcal{P}}$; satisfiability implies consistency.
\end{proof}

\subsection{Counter-Models vs Separation Formulas}
\label{subsec:cm-sep}

\begin{proposition}[Complementarity]
\label{prop:cm-sep}
The internal counter-model and the separation formula of Paper~I are
complementary tools:

\smallskip
\begin{center}
\begin{tabular}{lll}
\toprule
Property & Counter-Model & Separation Formula \\
\midrule
Scope & Single path $\mathcal{P}$ & Two paths $\mathcal{P}_1$,
  $\mathcal{P}_2$ \\
Establishes & $\mathcal{P} \nvdash \varphi$ &
  $\mathcal{P}_1 \vdash \varphi$, $\mathcal{P}_2 \vdash \neg\varphi$ \\
Primary use & Diagnostic + Prescription & Comparative analysis \\
\bottomrule
\end{tabular}
\end{center}

\smallskip
Neither subsumes the other; both are necessary for a complete account
of derivability in the Foundational Space.
\end{proposition}

\begin{definition}[Canonical Counter-Model]
\label{def:canonical-cm}
A counter-model $\mathcal{M}$ for $\varphi$ relative to $\mathcal{P}$
is \emph{canonical} if it is both minimal (no proper sub-model is
also a counter-model) and maximally informative (the structural
deficit $S(\mathcal{M},\varphi)$ is minimal among all counter-models
for $\varphi$ relative to $\mathcal{P}$).
\end{definition}

\begin{corollary}[Pre-Verification of Additive Correction]
\label{cor:cm-preverify}
An Additive Correction prescribed by a canonical internal
counter-model is automatically Minimal-Injection-compliant: every
prescribed injection is non-derivable from prior content, and the
set of injections is minimal.
\end{corollary}

% ============================================================
\section{The Machine-Readable Block Record}
\label{sec:record}
% ============================================================

\subsection{The Block Record Schema}
\label{subsec:schema}

\begin{definition}[Block Summary Record]
\label{def:record}
The \emph{Block Summary Record} of Block $\beta$ is a structured
document with the following fields:

\begin{lstlisting}[language=yaml, caption={Block Summary Record Schema}]
# Group I: Identity
ID:      TYPE-BLOCK-SEQ
Title:   [Short descriptive name]
Status:  Draft | Immutable

# Group II: Architectural Character
EpistemicStatus:     LN | FN | AN | MC
EpistemicJustification: [text or reference]
Conservativity:      Conservative | NonConservative
ConservativityBasis: [text or reference]

# Group III: Dependency and Content
Depends_on:  [TYPE-BLOCK-SEQ, ...]
Introduces:  [VOC-BLOCK-SEQ, ...]
Exports:     [TYPE-BLOCK-SEQ, ...]
Injects:     [AXM-BLOCK-SEQ, ...]
Derives:     [DEF/THM-BLOCK-SEQ, ...]

# Group IV: Decision and Correction
DecisionPoint:
  ID:                DP-BLOCK-SEQ | null
  Options:           [...] | null
  Selection:         [...] | null
  SelectionRationale: [text] | null
  IndependenceBasis: Objective | Interpretive | null
  UnChosenOptions:   [...] | null
ErratumStatus:
  ID:              ERR-SEQ | null
  Nature:          [...] | null
  CorrectionMode:  Additive | Branch | null
  CorrectionRef:   [...] | null

# Group V: Formal Payload
FormalPayload:
  Language:          Lean4 | Coq | Agda | Isabelle | Informal
  Source:            file:path | inline | null
  VerificationStatus: Compiled | ProofChecked | Reviewed | Unverified
\end{lstlisting}
\end{definition}

\begin{axiom}[Mandatory Block Record]
\label{ax:record}
Every Block $\beta$ in a \BOMA{}-compliant development must have a
Block Summary Record satisfying:
\begin{enumerate}
  \item \textbf{Completeness:} all fields populated;
  \item \textbf{Accuracy:} every field accurately reflects the
        Block's structural situation at introduction;
  \item \textbf{Contemporaneity:} produced at Block introduction,
        not retroactively;
  \item \textbf{Consistency:} structural declarations consistent
        with the formal payload.
\end{enumerate}
\end{axiom}

\subsection{Block Records and BOMA-Compliance}
\label{subsec:record-compliance}

\begin{proposition}[Block Records Establish Compliance]
\label{prop:record-compliance}
A \BOMA{}-compliant development in which every Block has a complete
and accurate Block Summary Record automatically satisfies all six
compliance properties of Definition~\ref{def:compliant}.
\end{proposition}

\begin{proof}
\textbf{P1:} \texttt{Depends\_on} and \texttt{Injects} enable
provenance tracing.
\textbf{P2:} \texttt{Epistemic\-Status}, \texttt{Conservativity}, and
all legibility fields implement Transparency.
\textbf{P3:} The record supports all four auditing operations;
\texttt{Verification\-Status} records mechanical verification.
\textbf{P4:} The separation of \texttt{Injects} and \texttt{Derives}
enforces Minimal Injection; Vocabulary-Axiom Decoupling requires
\texttt{Introduces} and \texttt{Injects} to appear in separate Blocks.
\textbf{P5:} \texttt{Decision\-Point} fields provide Branchability
documentation.
\textbf{P6:} \texttt{Status} enforces immutability;
\texttt{Erratum\-Status} implements the ERRATA mechanism.
\end{proof}

\subsection{Automated Compliance Checking}
\label{subsec:checker}

\begin{definition}[Compliance Checker]
\label{def:checker}
A \emph{Compliance Checker} verifies the following conditions from
the Block Records:
\begin{description}
  \item[C1] Every Block has a complete Record.
  \item[C2] The \texttt{Depends\_on} graph is acyclic.
  \item[C3] \texttt{Conservativity} is consistent with
            \texttt{Injects}.
  \item[C4] \texttt{EpistemicStatus} is consistent with
            \texttt{Conservativity} (Proposition~\ref{prop:cons-es}).
  \item[C5] Every $\MC$ Block has a fully populated
            \texttt{DecisionPoint} group.
  \item[C6] Every non-null \texttt{ErratumStatus.ID} refers to an
            existing ERRATA entry.
  \item[C7] \texttt{Exports} is the union of \texttt{Introduces},
            \texttt{Injects}, \texttt{Derives}; all IDs are unique;
            every \texttt{Depends\_on} ID is in \texttt{Exports} of
            a prior Block.
  \item[C8] No \texttt{Immutable} Block has any field other than
            \texttt{ErratumStatus} modified after the status was
            assigned.
\end{description}
All checks C1--C8 are decidable from a finite set of Block Records.
\end{definition}

% ============================================================
\section{Worked Example: The Pre-Numeric Kernel}
\label{sec:example}
% ============================================================

\subsection{Overview}
\label{subsec:ex-overview}

The Pre-Numeric Kernel of the Transparent Mathematics project spans
sixteen Bricks organized into two groups: the Zero-Group Constitution
(Bricks 000--006), containing meta-mathematical rules governing the
development's conduct; and the Pre-Numeric Kernel Proper (Bricks
007--015), containing formal mathematical content formalized in
Lean~4~\cite{Lean4,deMouraEtAl2015}.

The project's organizational discipline --- atomic Bricks,
vocabulary-axiom decoupling, explicit epistemic declarations,
immutability, formal ERRATA logging --- instantiates the six
compliance properties independently of \BOMA{}. The convergence
between the project's structure and \BOMA{}'s requirements is
evidence of the framework's descriptive adequacy.

\subsection{Compliance Assessment}
\label{subsec:ex-compliance}

\paragraph{P1 (Traceability).}
Substantially satisfied. The Global Unique Identifier system
($\mathrm{TYPE\text{-}BLOCK\text{-}SEQ}$) enables provenance tracing
for every exported entity. The dependency graph of the Lean~4
identifiers is tracked by the proof assistant itself.

\paragraph{P2 (Transparency).}
Largely satisfied. The \texttt{Necessity Type} declarations in the
project's machine-readable summaries map exactly onto the four-way
Epistemic Status taxonomy, and \texttt{Conservativity} declarations
are present in all Bricks 007--015. Minor gap: Bricks 000--006 use an
older format lacking the full Block Record schema of
Definition~\ref{def:record}.

\paragraph{P3 (Auditability).}
Mechanical Auditability achieved for Bricks 007--015 (Lean~4
compilation and proof checking). Interpretive Auditability for Bricks
000--006. The Compliance Checker of Definition~\ref{def:checker} has
not yet been implemented.

\paragraph{P4 (Minimal-Brick Construction).}
Fully satisfied. Vocabulary-Axiom Decoupling is rigorously maintained:
every vocabulary item (Bricks 008, 010, 011, 013) is introduced in a
dedicated Brick before any axiom governs it. Single-Purpose Commitment
is observed throughout.

\paragraph{P5 (Branchability).}
Satisfied. With both Brick~007 (logic choice) and Brick~011
(primitive positive distinction) classified as $\MC$:
\[
  \BI{\mathcal{P}_{\mathrm{PNK}}} = 2.
\]
Both Decision Points are documented with options, selections,
rationales, and independence bases.

\paragraph{P6 (Correctability).}
Satisfied through ERR-001's Additive Correction. No in-place
modification was made; the flaw is formally logged with a canonical
counter-model witness; the correction prescription is specified.

\subsection{The Architectural DAG}
\label{subsec:ex-dag}

The full Architectural DAG of the Pre-Numeric Kernel (Bricks 007--015)
has the following structure:

\begin{center}
\begin{small}
\begin{verbatim}
beta_007 (MC, NC) -- Logical Kernel
    |
    v
beta_008 (FN, NC) -- Domain E
    |
    +----------+----------+----------+
    v          v          v          v
beta_009    beta_010   beta_011   beta_013
(FN,NC)    (FN,NC)    (MC,NC)    (FN,NC)
Inhabit.   Anchor     Dist       Next
    |          |          |          |
    +----+-----+          +----+-----+
         v                    v
      beta_012            beta_014
      (FN,NC)             (FN,NC)
      Diversity           Stability
           |                   |
           +--------+----------+
                    v
                beta_015 (LN, C)
                Derived Predicate
\end{verbatim}
\end{small}
\end{center}

\noindent
Structural properties: 9~vertices, 13~edges, maximum in-degree~4
($\beta_{014}$), maximum out-degree~4 ($\beta_{008}$), 1~Conservative
node ($\beta_{015}$), 1~Convergence Point (candidate: $D(x)$
predicate in future Branches).

\begin{remark}
The DAG is non-linear at Brick~008 (out-degree~4) and Brick~014
(in-degree~4) as a direct consequence of Vocabulary-Axiom Decoupling:
the domain $\mathbf{E}$ is introduced once and used by four subsequent
Bricks (Proposition~\ref{prop:cons-es}).
\end{remark}

\subsection{ERR-001: Full \BOMA{} Analysis}
\label{subsec:ex-err}

\paragraph{Flaw Classification.}
$\mathrm{OverClaim}$: the Milestone-001 documentation asserted infinite
generative capacity, but this is not in
$\Th{\mathcal{P}_{\mathrm{PNK}}}$.

\paragraph{Correction Mode (Theorem~\ref{thm:mode}).}
$\mathrm{OverClaim}$ always requires Additive Correction regardless of
Epistemic Status. $\checkmark$

\paragraph{Counter-Model Analysis.}
The counter-model $\mathcal{M}_{\mathrm{ERR}}$:
\[
  \mathbf{E} = \{\bullet\},\quad
  o = \bullet,\quad
  \mathrm{Dist}(x,y) = \top,\quad
  \mathrm{Next}(x) = x.
\]
Compatibility: all axioms AXM-009-001, AXM-012-001, AXM-014-001
are satisfied. $\checkmark$\\
Falsification: no infinite sequence of pairwise distinct elements
exists in a one-element universe. $\checkmark$\\
Consistency confirmation (Proposition~\ref{prop:cm-lc}):
$\mathcal{P}_{\mathrm{PNK}}$ is locally consistent. $\checkmark$

\paragraph{Structural Deficit and Prescription.}
Two conditions permit $\mathcal{M}_{\mathrm{ERR}}$:
\begin{align*}
  &\mathrm{Dist}(\bullet,\bullet) = \top \quad
    \text{(reflexivity of Dist not ruled out)}, \\
  &\mathrm{Next}(\bullet) = \bullet \quad
    \text{(fixed-point of Next not ruled out)}.
\end{align*}
Minimal Additive Correction (Theorem~\ref{thm:cm-prescribe}):
\begin{align*}
  \psi_{\mathrm{irr}} &:
    \forall x : \mathbf{E}.\; \neg\mathrm{Dist}(x,x), \\
  \psi_{\mathrm{inj}} &:
    \forall x, y : \mathbf{E}.\;
    \mathrm{Next}(x) = \mathrm{Next}(y) \to x = y.
\end{align*}

\paragraph{Canonical Character.}
$\mathcal{M}_{\mathrm{ERR}}$ is minimal (one-element domain, no
proper sub-models) and maximally informative (identifies exactly two
structural deficits). By Definition~\ref{def:canonical-cm}, it is a
canonical counter-model. $\checkmark$

\paragraph{Pre-Verification (Corollary~\ref{cor:cm-preverify}).}
The canonical counter-model certifies Minimal-Injection compliance of
the Additive Correction before any new Brick is introduced. $\checkmark$

\subsection{Block Records: Three Representative Examples}
\label{subsec:ex-records}

\begin{example}[Brick\_007 --- Methodological Choice, Non-Conservative]
\label{ex:brick007}
\begin{lstlisting}[language=yaml]
ID:      AXM-007-LOG
Title:   Intuitionistic Logical Kernel
Status:  Immutable

EpistemicStatus:      MethodologicalChoice
EpistemicJustification: >
  Classical first-order logic was a genuine alternative.
  Intuitionistic logic was selected to require explicit
  construction for all existential claims. Independence basis:
  objective (Goedel-Gentzen negative translation).
Conservativity:      NonConservative
ConservativityBasis: >
  Introduces AXM-007-001 and AXM-007-002, neither derivable
  from prior content (there is no prior content).

Depends_on:  []
Introduces:  []
Exports:     [AXM-007-001, AXM-007-002]
Injects:     [AXM-007-001, AXM-007-002]
Derives:     []

DecisionPoint:
  ID:                DP-007-001
  Options:           [Classical FOL, Intuitionistic MSFOL, Relevant]
  Selection:         Intuitionistic MSFOL
  SelectionRationale: >
    Preserves constructive existence; prevents non-constructive
    existence proofs from infiltrating the kernel.
  IndependenceBasis: Objective
  UnChosenOptions:   [Classical FOL, Relevant logic]

ErratumStatus:
  ID: null

FormalPayload:
  Language:           Lean4
  Source:             Formalization/PreNumericKernel.lean
  VerificationStatus: Compiled
\end{lstlisting}
\end{example}

\begin{example}[Brick\_012 --- Foundational Necessity, Non-Conservative]
\label{ex:brick012}
\begin{lstlisting}[language=yaml]
ID:      AXM-012-001
Title:   Anchor-Relative Diversity Axiom
Status:  Immutable

EpistemicStatus:      FoundationalNecessity
EpistemicJustification: >
  Given the development's goal of structural diversity, at least
  two distinguishable elements are needed. Anchor-Relative
  formulation is minimal; no weaker alternative achieves the
  structural goal.
Conservativity:      NonConservative
ConservativityBasis: >
  AXM-012-001 not derivable from prior vocabulary and inhabitation
  axioms. Consistency witness: two-element model {o, a} with
  Dist(a,o) = True.

Depends_on:  [VOC-008-001, AXM-009-001, VOC-010-001, VOC-011-001]
Introduces:  []
Exports:     [AXM-012-001]
Injects:     [AXM-012-001]
Derives:     []

DecisionPoint:
  ID: null

ErratumStatus:
  ID: null

FormalPayload:
  Language:           Lean4
  Source:             Formalization/PreNumericKernel.lean
  VerificationStatus: Compiled
\end{lstlisting}
\end{example}

\begin{example}[Brick\_015 --- Logical Necessity, Conservative]
\label{ex:brick015}
\begin{lstlisting}[language=yaml]
ID:      DEF-015-001
Title:   Derived Predicate and Generative Consequences
Status:  Immutable

EpistemicStatus:      LogicalNecessity
EpistemicJustification: >
  DEF-015-001 is a pure macro (eliminable). THM-015-001,
  THM-015-002, THM-015-003 are strict logical consequences
  of AXM-012-001 and AXM-014-001, verified by Lean4.
Conservativity:      Conservative
ConservativityBasis: >
  No new axioms or vocabulary. All content derivable from
  prior Bricks. Lean4 proofs require no additional axioms.

Depends_on:  [VOC-011-001, AXM-012-001, VOC-013-001, AXM-014-001]
Introduces:  []
Exports:     [DEF-015-001, THM-015-001, THM-015-002, THM-015-003]
Injects:     []
Derives:     [DEF-015-001, THM-015-001, THM-015-002, THM-015-003]

DecisionPoint:
  ID: null

ErratumStatus:
  ID:             ERR-001
  Nature:         OverClaim
  CorrectionMode: Additive
  CorrectionRef:  [AXM-MILESTONE002-IRREF, AXM-MILESTONE002-INJ]

FormalPayload:
  Language:           Lean4
  Source:             Formalization/PreNumericKernel.lean
  VerificationStatus: ProofChecked
\end{lstlisting}
\end{example}

\subsection{Architectural Insights from the Example}
\label{subsec:ex-insights}

The Pre-Numeric Kernel analysis generates five architectural
observations that extend the theoretical development of preceding
sections.

\begin{observation}[Vocabulary Bricks Are Systematically
Non-Conservative]
Every Brick introducing a vocabulary item (Bricks 008, 010, 011, 013)
is Non-Conservative: vocabulary introduction expands the signature,
enabling new formulas even before any axiom governs the new symbol.
This is a general pattern: \emph{vocabulary introduction is always
Non-Conservative}.
\end{observation}

\begin{observation}[The First Conservative Threshold]
Brick~015 is the first Conservative Brick in the object-level
development, appearing after 8~Non-Conservative object-level Bricks.
This \emph{First Conservative Threshold} marks the moment at which
the development has accumulated sufficient foundational content to
generate non-trivial derived consequences.
\end{observation}

\begin{observation}[Counter-Models Are More Informative Than
Non-Derivability Proofs]
An abstract proof that $\varphi_\infty$ is not derivable from
$\mathcal{P}_{\mathrm{PNK}}$ would establish the over-claim but
would not specify the correction. The canonical counter-model
$\mathcal{M}_{\mathrm{ERR}}$ simultaneously establishes the
non-derivability \emph{and} prescribes the minimal Additive Correction
--- a general consequence of Theorem~\ref{thm:cm-prescribe}.
\end{observation}

\begin{observation}[The Architectural DAG Is Irreducible]
The Architectural DAG of the Pre-Numeric Kernel has 9~vertices and
13~edges. Removing any vertex breaks the connectivity required for
Brick~015 to derive its theorems; removing any edge removes a
dependency genuinely required by the dependent Brick. The DAG is
irreducible --- a structural consequence of \BOMA{}-compliance:
\emph{every \BOMA{}-compliant development satisfying Minimal-Brick
Construction produces an irreducible Architectural DAG}.
\end{observation}

\begin{observation}[The Corrected Development Prescribes Milestone~002]
The canonical counter-model prescribes exactly two new Non-Conservative
Bricks for Milestone~002:
\begin{align*}
  \text{Brick\_016} &:
    \forall x : \mathbf{E}.\;\neg\mathrm{Dist}(x,x)
    \quad \text{(Irreflexivity, } \ES{} = \FN\text{)}, \\
  \text{Brick\_017} &:
    \forall x,y : \mathbf{E}.\;
    \mathrm{Next}(x) = \mathrm{Next}(y) \to x = y
    \quad \text{(Injectivity, } \ES{} = \FN\text{)}.
\end{align*}
These two injections are independently necessary (neither derivable
from the other plus prior content) and jointly sufficient to rule out
$\mathcal{M}_{\mathrm{ERR}}$.
\end{observation}

% ============================================================
\section{Relation to \BOMA{} Paper~I}
\label{sec:paper1}
% ============================================================

\subsection{What This Paper Extends}
\label{subsec:rel-extends}

This paper extends Paper~I in four directions, none contradicting its
results.

\textbf{Extension~1 --- Operational Specification.}
Paper~I defined the structural components; this paper defines the
operational properties those components jointly produce when correctly
applied.

\textbf{Extension~2 --- Finer Epistemic Granularity.}
The binary Forced Step / Genuine Choice of Paper~I is replaced by the
four-way Epistemic Status taxonomy.

\textbf{Extension~3 --- Internal Diagnostic Tools.}
Paper~I's tools (separation formulas, structural comparison) operate
between paths. The Internal Counter-Model operates within a single
path, complementing those tools.

\textbf{Extension~4 --- Correction Procedures.}
Paper~I established Local Consistency as the validity criterion but
did not address correction procedures for flaws within a path. This
paper provides those procedures through the two-mode framework of
\S\ref{sec:correction}.

\subsection{Formal Relationships}
\label{subsec:rel-formal}

\begin{proposition}[Refinement]
\label{prop:refinement}
All structural results of Paper~I remain valid under the four-way
Epistemic Status taxonomy (Proposition~\ref{prop:es-compat}).
\BOMA{}-compliance strictly strengthens Architectural Transparency
(Proposition~\ref{prop:comp-at}). The Architectural DAG extends the
Branch Tree as its Decision-Level Projection
(Definition~\ref{def:projection}). The Internal Counter-Model and
Separation Formulas are complementary, non-redundant tools
(Proposition~\ref{prop:cm-sep}).
\end{proposition}

\subsection{Minimal Modifications to Paper~I}
\label{subsec:rel-mods}

Two additions to Paper~I are indicated by the results of this paper:

\textbf{Addition~1.} At the end of Paper~I, Remark~4.3.11, add:
\begin{quote}
\textit{A complete four-way classification of the sources of necessity
is provided in Paper~II of this series, \S3, where the binary
distinction of the present paper is extended to a four-way Epistemic
Status taxonomy. The binary classification of Forced Step and Genuine
Choice is preserved as a coarsening of that taxonomy.}
\end{quote}

\textbf{Addition~2.} After Paper~I, Definition~4.4.12, add
Remark~4.4.12a (as introduced in \S\ref{sec:dag}):
\begin{quote}
\textit{At the level of individual Blocks, the full architectural
structure may be richer than a tree: a Block may be required in
multiple branches without local modification, constituting a
Convergence Point. The formal treatment of Convergent Blocks and the
full Architectural DAG is developed in Paper~II, \S6.}
\end{quote}

\subsection{What Remains for Paper~III}
\label{subsec:rel-paper3}

Papers~I and~II together establish \BOMA{} as a theoretical framework
(Paper~I) with a complete operational specification (Paper~II).
Paper~III will provide a fully developed mathematical example:

\begin{enumerate}
  \item \textbf{Milestone~002 Completion}: introduce Bricks~016
        (irreflexivity) and~017 (injectivity) as the prescribed
        Additive Correction;
  \item \textbf{Construction of $\mathbb{N}$}: formal development
        within \BOMA{}, building on the corrected Pre-Numeric Kernel;
  \item \textbf{First Branch Development}: open a Branch at the
        Brick~011 Decision Point and identify the first non-trivial
        Convergent Blocks between the two paths;
  \item \textbf{Architectural Metrics}: track the growth of the DAG,
        Branchability Index, Conservativity Ratio, and cumulative
        injection weight as the development matures.
\end{enumerate}

% ============================================================
\section{Conclusion}
\label{sec:conclusion}
% ============================================================

\subsection{What This Paper Has Not Done}
\label{subsec:conc-not}

This paper has not proposed a new mathematical theory, revised
Paper~I, or completed the Transparent Mathematics project's
Milestone~002. The Additive Correction prescribed by ERR-001 has been
specified and analyzed, but not introduced as formal Bricks. The
analysis of \S\ref{sec:example} is diagnostic, not constructive.

\subsection{What This Paper Has Done}
\label{subsec:conc-done}

We have formalized \BOMA{}-compliance as a conjunction of six
properties (Definition~\ref{def:compliant}), enriched the epistemic
classification to four-way status (\S\ref{sec:epistemic}), introduced
the mandatory Conservativity Declaration (\S\ref{sec:cons}), extended
the Branch Tree to the Architectural DAG and introduced Convergent
Blocks as the architectural criterion for foundational invariance
(\S\ref{sec:dag}), distinguished two modes of correction
(\S\ref{sec:correction}), formalized the Internal Counter-Model as a
structural prescription (\S\ref{sec:counter}), specified the
Machine-Readable Block Record (\S\ref{sec:record}), and demonstrated
all contributions through a complete worked example
(\S\ref{sec:example}).

\subsection{The Two Central Claims}
\label{subsec:conc-claims}

\textbf{First claim:} \BOMA{}-compliance is not merely a theoretical
ideal but a practicable discipline. The Pre-Numeric Kernel satisfies
five of the six compliance properties fully and the sixth
substantially, through an independently motivated organizational
discipline that converged on \BOMA{}-compatible structure without
prior knowledge of the framework.

\textbf{Second claim:} The Convergent Block is the precise
architectural criterion for foundational invariance. A mathematical
result that appears identically in two paths diverging at a Decision
Point is path-independent with respect to that Decision Point ---
invariant under the foundational choice made there. This is the
formal expression of ``foundationally neutral'' mathematics: not a
philosophical characterization but a structural certificate,
computable from the Architectural DAG, establishing precisely which
results survive which foundational choices.

\subsection{A Remark on Mathematical Organization}
\label{subsec:conc-remark}

This paper has argued throughout that how mathematics is organized is
not epistemically neutral. The organizational paradigm of a
mathematical development determines what is visible, traceable,
correctable, and comparable. The Transparent Mathematics project
demonstrates that a development can be simultaneously mathematically
rigorous (mechanically verified in Lean~4) and architecturally
transparent (Bricks documented, Epistemic Status declared, ERRATA
maintained). The compatibility requires discipline, infrastructure,
and a commitment to recording the construction history alongside the
mathematical content. But it is achievable --- from the very
beginning of a foundational development.

The architecture is open. The Convergent Blocks are waiting to be
discovered. The Branches are waiting to be developed. And the
mathematical territory beyond the Pre-Numeric Kernel is structurally
present, determinately reachable, and waiting to be made actual by
the Bricks that will traverse it.

% ============================================================
% BIBLIOGRAPHY
% ============================================================

\newpage
\printbibliography

% ============================================================
\end{document}
% ============================================================
